% =====================================================================
\needspace{0.62\textheight}
\appendix
\renewcommand{\thesection}{\Alph{section}}
\section{Formulario del Tema 3}
% =====================================================================
\begingroup
\small
\renewcommand{\arraystretch}{1.85}
\newcommand{\grupo}[1]{\multicolumn{3}{l}{\color{acento}\espaciado[200]{#1}}\\[-2pt]}
\noindent
\begin{tabularx}{\textwidth}{@{}l>{\centering\arraybackslash}Xr@{}}
\toprule
Magnitud o resultado & Expresión & Referencia\\
\midrule
\grupo{Máquinas térmicas}
Motor & $\eta=-W/Q_h=1+Q_c/Q_h$;\ varios tramos: $|W|/Q_{\mathrm{abs}}$ & ec.~\eqref{eq:rendimientos}\\
Refrigerador y bomba & $\varepsilon=Q_c/W,\quad\nu=-Q_h/W=\varepsilon+1$ & ec.~\eqref{eq:rendimientos}\\
Carnot & $\eta_C=1-\dfrac{T_c}{T_h},\quad\varepsilon_C=\dfrac{T_c}{T_h-T_c},\quad\nu_C=\dfrac{T_h}{T_h-T_c}$ & ec.~\eqref{eq:carnot}\\
Otto y Brayton & $\eta=1-r^{1-\gamma}$,\quad $\eta=1-r_p^{(1-\gamma)/\gamma}$ & ec.~\eqref{eq:etaotto}\\
\midrule
\grupo{Entropía}
Definición & $\dd S=\dbar Q_{\mathrm{rev}}/T$ & ec.~\eqref{eq:dS}\\
Balance & $\Delta S=\int\dbar Q/T_F+S_{\mathrm{gen}},\quad S_{\mathrm{gen}}\geq0$ & ec.~\eqref{eq:balance}\\
Universo & $\Delta S_{\mathrm{univ}}=\Delta S_{\mathrm{sis}}+\Delta S_{\mathrm{ent}}\geq0$ & ec.~\eqref{eq:universo}\\
Gas ideal & $\Delta S=C_V\ln\frac{T_2}{T_1}+nR\ln\frac{V_2}{V_1}=C_p\ln\frac{T_2}{T_1}-nR\ln\frac{p_2}{p_1}$ & ec.~\eqref{eq:Sgas}\\
Cuerpo y foco & $\Delta S=C\ln(T_f/T_i)$;\quad $\Delta S_F=-Q_{\mathrm{sis}}/T_F$ & tabla~\ref{tab:dS}\\
\midrule
\grupo{Ecuaciones prácticas}
Relación fundamental & $\dd U=T\,\dd S-p\,\dd V$ & ec.~\eqref{eq:fundamental}\\
Ec.\ de la energía & $\pdi UVT=T\pdi pTV-p$ & ec.~\eqref{eq:energetica}\\
Ecuaciones $T\,\dd S$ & $\dd S=\frac{C_V}{T}\dd T+\pdi pTV\dd V=\frac{C_p}{T}\dd T-\pdi VTp\dd p$ & ec.~\eqref{eq:TdS}\\
$C_p-C_V$ general & $TV\alpha^2/\kappa_T\geq0$ & ec.~\eqref{eq:CpCVgeneral}\\
\midrule
\grupo{Trabajo útil}
Trabajo perdido & $W_{\mathrm{perdido}}=T_cS_{\mathrm{gen}}$ & ec.~\eqref{eq:perdido}\\
Cuerpos finitos & $\Delta S_{\mathrm{univ}}=0\Rightarrow T_f=\sqrt{T_1T_2}$ (iguales) & probl.~38\\
\bottomrule
\end{tabularx}
\endgroup

% =====================================================================
\begingroup
\titleformat{\section}{\Large\bfseries\color{navy}}{}{0em}{}[\vspace{2pt}{\color{navy}\titlerule[0.9pt]}]
\renewcommand{\refname}{Referencias}
\makeatletter
\renewenvironment{thebibliography}[1]
  {\section*{\refname}\addcontentsline{toc}{section}{\refname}\markboth{\refname}{}%
   \list{\@biblabel{\@arabic\c@enumiv}}{\settowidth\labelwidth{\@biblabel{#1}}\leftmargin\labelwidth
     \advance\leftmargin\labelsep\itemsep 4pt\usecounter{enumiv}\let\p@enumiv\@empty
     \renewcommand\theenumiv{\@arabic\c@enumiv}}\sloppy}
  {\def\@noitemerr{\@latex@warning{Empty `thebibliography' environment}}\endlist}
\makeatother
\clearpage
\begin{thebibliography}{9}
\bibitem{gonzalez1} J.~González Ayala. \emph{Termodinámica I}, capítulo 3: Segunda ley, y Lista de
  problemas \#3. Apuntes de la asignatura. Grado en Física, Universidad de Salamanca, 2026.
\bibitem{gonzalez2} J.~González Ayala. \emph{Tema 3. Segundo principio: entropía, reversibilidad y límites del
  trabajo}. Transparencias. Universidad de Salamanca, curso 2026--27.
\bibitem{parcial} \emph{Termodinámica I. Prueba de evaluación continua (temas 1--3)}. Universidad de
  Salamanca, 31 de octubre de 2025.
\bibitem{martin2} D.~Martín Guerrero. \emph{Primera ley (Termodinámica I, Tema 2)}. Apuntes formalizados
  con seminarios prácticos. Grado en Física, Universidad de Salamanca, 2026.
\bibitem{callen} H.~B.~Callen. \emph{Thermodynamics and an Introduction to Thermostatistics}, 2.ª ed. Wiley, 1985.
\bibitem{zemansky} M.~W.~Zemansky y R.~H.~Dittman. \emph{Calor y termodinámica}, 6.ª ed. McGraw-Hill, 1990.
\end{thebibliography}
\endgroup
